% Part: proof-theory
% Chapter: cut-elimination
% Section: ce-topmost

\documentclass[../../../include/open-logic-section]{subfiles}

\begin{document}

\olfileid{pt}{cut}{top}

\olsection{સૌથી ઉપરના કટ દૂર કરવાં}

\begin{defn}
\CutCS{} નિયમથી પૂરી થતી !!a{proof}~$\pi$ લો,
\[
\AxiomC{}
\RightLabel{$\pi_1$}
\Deduce$\Gamma \fCenter \Delta, !A$
\AxiomC{}
\RightLabel{$\pi_2$}
\Deduce$!A, \Gamma \fCenter \Delta$
\RightLabel{\CutCS}
\BinaryInf$\Gamma \fCenter \Delta$
\DisplayProof
\]
અને તેમાં બીજું કોઈ \CutCS{} અનુમાન ન હોય.
~$\pi$ની \emph{કટ-ઊંચાઈ}
$\cheight{\pi} = \pheight{\pi_1} + \pheight{\pi_2}$ છે.
~$\pi$નો \emph{કટ-ક્રમ}
$\cutrank{\pi} = \depth{!A}$ છે.
\end{defn}

\begin{lem}\ollabel{lem:cut-adm-G3c}
\CutCS{} નિયમ~\Log{G3c}માં સ્વીકાર્ય છે.
એટલે, જો~$\pi$ એ માત્ર એક~\CutCS{}થી
પૂરી થતી અને અન્યથા માત્ર~$\Log{G3c}$ના નિયમો
વાપરતી !!a{proof} હોય, તો એ જ અંતિમ સિક્વન્ટની
!!a{proof}~$\pi'$ એ~\Log{G3c}માં છે.
\end{lem}

\begin{proof}
સાબિતી કટની !!{height} અને ક્રમ પર
દ્વિ-આગમનથી છે. !!{proof}~$\pi$નો અંત આ રીતે થાય છે:
\[
\AxiomC{}
\RightLabel{$\pi_1$}
\Deduce$\Gamma \fCenter \Delta, !A$
\AxiomC{}
\RightLabel{$\pi_2$}
\Deduce$!A, \Gamma \fCenter \Delta$
\RightLabel{\CutCS}
\BinaryInf$\Gamma \fCenter \Delta$
\DisplayProof
\]

આગમનનો આધાર: $\cheight{\pi} = 0$ અને
$\cutrank{\pi} = 0$. કારણ કે
$\cheight{\pi} = \pheight{\pi_1} + \pheight{\pi_2} = 0$,
$\Gamma \Sequent \Delta, !A$ અને
$!A, \Gamma \Sequent \Delta$ બંને સ્વયંસિદ્ધ છે.
બે કિસ્સા છે: (a)~તેમાંથી કોઈ એકમાં~$!A$
મુખ્ય નથી, અથવા (b)~બંનેમાં~$!A$ મુખ્ય છે.
નીચે બંને કિસ્સામાં
$\Gamma \Sequent \Delta$ સ્વયંસિદ્ધ છે
એ બતાવીશું (અને તે માટે આગમન-પૂર્વધારણા
જરૂરી નથી).

હવે~$\pi_1$ અને $\pi_2$ સ્વયંસિદ્ધ છે કે
અનુમાનથી પૂરી થાય છે, અને એ અનુમાનમાં
$!A$ મુખ્ય છે કે નહીં, તેના આધારે
કિસ્સા જુદા પાડીએ. કિસ્સા~A અને~B
આગમનનો આધાર આવરી લે છે.
આગમન-પગલામાં
$\cheight{\pi} > 0$ અથવા
$\cutrank{\pi} > 0$ (અથવા બંને)
ધારીએ છીએ. આગમન-પૂર્વધારણા મુજબ,
માત્ર એક છેલ્લા~\CutCS{}થી પૂરી થતી
એવી કોઈ પણ !!{proof}માંથી કટ દૂર કરી શકાય,
જેનો કટ-ક્રમ $= \cutrank{\pi}$ અને
કટ-!!{height} $< \cheight{\pi}$ હોય,
અથવા કટ-ક્રમ~$< \cutrank{\pi}$ હોય.
જ્યારે $\cheight{\pi} = 0$ હોય
(બંને આધારવિધાન સ્વયંસિદ્ધ હોય), ત્યારે
કિસ્સા~A અને~B લાગુ પડે છે.
$\cheight{\pi} > 0$ના કિસ્સા C--Dમાં આવે છે.

\paragraph{કિસ્સો A:}
કોઈ~$!B$ એ~$\Gamma$ અને~$\Delta$ બંનેમાં છે.
તો~$\Gamma \Sequent \Delta$ સ્વયંસિદ્ધ છે
(જો~$!B$ પરમાણ્વીય હોય), અથવા
\olref[seq][ex]{prop:G3c-axioms} મુજબ
\Log{G3c}માં તેની \CutCS{} વિનાની
!!a{proof}~$\pi'$ છે. આ આગમનના આધારનો
કિસ્સો~(a) આવરી લે છે: જો
$\Gamma \Sequent \Delta, !A$ કે
$!A, \Gamma \Sequent \Delta$ સ્વયંસિદ્ધ હોય,
પણ તેમાં~$!A$ મુખ્ય ન હોય,
તો કોઈ બીજું પરમાણ્વીય~$!B$
એ~$\Gamma$ અને~$\Delta$ બંનેમાં હોવું જોઈએ.
જો કોઈ એક આધારવિધાન એવું સ્વયંસિદ્ધ હોય
જેમાં~$!A$ મુખ્ય નથી,
તો બીજું આધારવિધાન નિયમનો નિષ્કર્ષ હોય
ત્યારે પણ આ આગમન-પગલું આવરી લે છે.

\paragraph{કિસ્સો B:}
બંને આધારવિધાન સ્વયંસિદ્ધ છે અને
$!A$ મુખ્ય છે, તેથી પરમાણ્વીય છે.
ડાબું આધારવિધાન
$\Gamma \Sequent \Delta, !A$ લો.
$!A$ મુખ્ય હોવાથી તે~$\Gamma$માં હોવું જોઈએ
(નહીં તો~$\Gamma \Sequent \Delta, !A$
સ્વયંસિદ્ધ ન હોત). એ જ રીતે
$!A$ એ~$\Delta$માં હોવું જોઈએ,
નહીં તો~$!A, \Gamma \Sequent \Delta$
સ્વયંસિદ્ધ ન હોત. તેથી પરમાણ્વીય~$!A$
એ~$\Gamma$ અને~$\Delta$ બંનેમાં છે,
એટલે~$\Gamma \Sequent \Delta$
સ્વયંસિદ્ધ છે. આ આગમનના આધારનો
કિસ્સો~(b) આવરી લે છે.

\begin{center}
કિસ્સા A અને Bની વૈકલ્પિક રજૂઆત
\end{center}

\paragraph{કિસ્સો A:}
કટનાં બે આધારવિધાનમાંથી એક,
પરમાણ્વીય~$!C$ સાથે
$!C, \Gamma' \Sequent \Delta', !C$
સ્વરૂપનું સ્વયંસિદ્ધ છે.
જો~$!C$ એ કટ !!{formula}~$!A$ ન હોય,
તો~$!C$ એ~$\Gamma$ અને~$\Delta$ બંનેમાં
હોવું જોઈએ. ત્યારે
$\Gamma \Sequent \Delta$ સ્વયંસિદ્ધ છે
અને તેને~$\pi'$ લઈ શકાય.
જો~$!C \ident !A$ એ કટ !!{formula} હોય,
તો બે શક્યતાઓ છે: ડાબું આધારવિધાન સ્વયંસિદ્ધ
હોય તો $!A \in \Gamma$ અને
$\Gamma = \Gamma', !A$;
જમણું સ્વયંસિદ્ધ હોય તો
$!A \in \Delta$ અને
$\Delta = \Delta', !A$.
પહેલા કિસ્સામાં~$\pi_2$ એ
$!A, !A, \Gamma' \Sequent \Delta$ની
!!a{proof} છે; \LeftR{\Contraction}ની
સ્વીકાર્યતા
(\olref[seq][inv]{prop:G3c-cont-adm}) વડે
$\Gamma \Sequent \Delta$ની !!a{proof} મળે છે.
બીજા કિસ્સામાં~$\pi_1$ એ
$\Gamma \Sequent \Delta', !A, !A$ની
!!a{proof} છે; \RightR{\Contraction}ની
સ્વીકાર્યતા વડે
$\Gamma \Sequent \Delta$ની !!a{proof} મળે છે.
\sourcecorrection{OLMD-200}{મૂળમાં બીજા કિસ્સામાં $\pi_2$ લખ્યું છે; બેવડા ઉત્તરાંગવાળી સાબિતી ડાબા આધારવિધાનની $\pi_1$ છે.}

\paragraph{કિસ્સો B:}
આધારવિધાનોમાંથી એક
$\lfalse, \Gamma' \Sequent \Delta'$
સ્વરૂપનું સ્વયંસિદ્ધ છે.
જો ડાબું આધારવિધાન આવું હોય,
તો $\lfalse \in \Gamma$ અને
$\Gamma \Sequent \Delta$ પણ સ્વયંસિદ્ધ છે.
જો જમણું આવું હોય, તો
$\lfalse \in \Gamma$ હોઈ શકે
(ત્યારે પણ $\Gamma \Sequent \Delta$
સ્વયંસિદ્ધ છે), અથવા
$!A \ident \lfalse$ હોઈ શકે.
પછીના કિસ્સામાં~$\pi_1$ એ
$\Gamma \Sequent \Delta, \lfalse$ની
!!a{proof} છે. ~$\pi_1$ના દરેક
સિક્વન્ટના ઉત્તરાંગમાંથી~$\lfalse$ની
એક નકલ કાઢી તેને
~$\Gamma \Sequent \Delta$ની
!!a{proof}માં ફેરવી શકાય.
(કસરત: $\pheight{\pi_1}$ પર
આગમનથી આ ચકાસો.)

હવે માનો કે કિસ્સો~A કે~B લાગુ પડતો નથી.
બાકી રહેલા કિસ્સાઓમાં
$\cheight{\pi} > 0$ છે, એટલે
ઓછામાં ઓછું એક આધારવિધાન
કોઈ અનુમાનનો નિષ્કર્ષ છે.

\paragraph{કિસ્સા C અને D: કટનું સ્થાનાંતર.}
કિસ્સા~C અને~Dની બધી શક્યતાઓમાં
એક સામાન્ય ઢબ છે. શરૂઆતમાં
~\CutCS{}થી પૂરી થતી !!a{proof} લઈએ.
તેનું એક આધારવિધાન નિયમ~$R$થી મળે છે,
જેમાં કટ !!{formula}~$!A$ મુખ્ય નથી
(બીજું કોઈ !!{formula}~$!D$ મુખ્ય છે).
આ !!{proof}માં \CutCS{} એ નિયમ~$R$
પછી આવે છે; યોજનાત્મક રીતે આ પ્રમાણે:
\[
\AxiomC{}
\RightLabel{$\pi_1$}
\Deduce$\Gamma \fCenter \Delta', !D, !A$
\AxiomC{}
\RightLabel{$\pi_3$}
\Deduce$!A, \Gamma, \Pi \fCenter \Delta', \Lambda$
\RightLabel{$R$}
\UnaryInf$!A, \Gamma \fCenter \Delta', !D$
\RightSubproofLabel{$\pi_2$}
\RightLabel{\CutCS}
\BinaryInf$\Gamma \fCenter \Delta', !D$
\DisplayProof
\]
આ માત્ર ત્યારેનું ચિત્ર છે જ્યારે~$R$
એક આધારવિધાનવાળો જમણો નિયમ હોય,
$A$ એ~$R$નું મુખ્ય !!{formula} ન હોય,
અને જે !!{formula}~$!D$ \emph{મુખ્ય}
છે તે~\CutCS{}ના જમણા આધારવિધાનના
ઉત્તરાંગમાં હોય. જો~$!D$ પૂર્વાંગમાં
હોય (એટલે~$R$ ડાબો નિયમ હોય)
અથવા~\CutCS{}ના ડાબા આધારવિધાનમાં
હોય, તો ચિત્ર અલગ હશે.
$\Pi$ અને~$\Lambda$નાં !!{formula}s
નિયમ~$R$નાં સક્રિય !!{formula}s દર્શાવે છે.

આ ક્રમ ઉલટાવી એવી !!a{proof} લઈએ
જેમાં~$R$ એ~\CutCS{} પછી આવે છે:
\[
\AxiomC{}
\RightLabel{$\pi_1'$}
\Deduce$\Gamma, \Pi \fCenter \Delta', !D, \Lambda, !A$
\AxiomC{}
\RightLabel{$\pi_3'$}
\Deduce$!A, \Gamma, \Pi \fCenter \Delta', !D, \Lambda$
\RightLabel{\CutCS}
\BinaryInf$\Gamma, \Pi \fCenter \Delta', !D, \Lambda$
\RightLabel{$R$}
\UnaryInf$\Gamma \fCenter \Delta', !D, !D$
\DisplayProof
\]
\sourcecorrection{OLMD-201}{મૂળ યોજનાત્મક વૃક્ષના અંતિમ સિક્વન્ટમાં $\Gamma$ પછી વધારાનો અલ્પવિરામ છે; પૂર્વાંગ માત્ર $\Gamma$ છે.}
!!{height} જાળવતા શિથિલીકરણથી
અનુક્રમે $\pi_1$ અને $\pi_3$માંથી
$\pi_1'$ અને $\pi_3'$ મેળવીએ છીએ.

\CutCS{} અને~$R$નો ક્રમ બદલી દીધો હોવાથી
તેને ``કટને~$R$ની ઉપર ખસેડવો'' કહે છે.
એથી~\CutCS{}ના જમણા આધારવિધાનમાં
પૂરી થતી !!{proof}ની ઊંચાઈ,
અને તેથી કટની !!{height}, ઘટે છે.
એટલે નવા~\CutCS{}માં પૂરી થતી
ઉપ-!!{proof}sની !!{height} અનુક્રમે
$\pi_1$ અને $\pi_3$થી વધારે નથી;
જમણી બાજુનું એક અનુમાન કાઢ્યું હોવાથી
કટની !!{height} ઘટે છે.
હવે આગમન-પૂર્વધારણા લાગુ પડે છે
અને~\CutCS{} અનુમાન દૂર થાય છે.
છેલ્લે~$\RightR{\Contraction}$ની
સ્વીકાર્યતા વડે~$!D$ની વધારાની નકલ
દૂર કરીએ છીએ.

\paragraph{કિસ્સો C:}
$\pi_1$ અને $\pi_2$માંથી ઓછામાં ઓછી એક
એવા એક-આધારવિધાનવાળા નિયમ~$R$થી પૂરી થાય છે
જેમાં~$!A$ મુખ્ય નથી.
કુલ 18 કિસ્સા છે: \CutCS{}ના ડાબા
કે જમણા આધારવિધાનમાં~$!A$ મુખ્ય નથી
અને નવ એક-આધારવિધાનવાળા નિયમોમાંથી
(\LeftR{\lnot}, \RightR{\lnot},
\RightR{\lor}, \LeftR{\land},
\RightR{\lif} તથા ચાર પરિમાણક નિયમો)
કયો~$R$ છે, તેના પરથી.
માત્ર બે ઉદાહરણ જોઈએ:
~$\pi_2$ના છેલ્લા અનુમાનમાં
$!A$ મુખ્ય નથી અને તેનો અંત
$\RightR{\lif}$ કે
\LeftR{\lexists}થી થાય છે.
બીજા કિસ્સા સમાન છે અને
કસરત તરીકે છોડીએ છીએ.

માનો કે~$\pi$નો અંત આ રીતે થાય છે:
\[
\AxiomC{}
\RightLabel{$\pi_1$}
\Deduce$\Gamma \fCenter \Delta', !B \lif !C, !A$
\AxiomC{}
\RightLabel{$\pi_3$}
\Deduce$!A, !B, \Gamma \fCenter \Delta', !C$
\RightLabel{$\RightR{\lif}$}
\UnaryInf$!A, \Gamma \fCenter \Delta', !B \lif !C$
\RightSubproofLabel{$\pi_2$}
\RightLabel{\CutCS}
\BinaryInf$\Gamma \fCenter \Delta', !B \lif !C$
\DisplayProof
\]
અહીં $\Delta = \Delta', !B \lif !C$.
~$\pi_1$ પર !!{height} જાળવતું
શિથિલીકરણ
(\olref[seq][adm]{prop:weak-G3c-adm}) લગાવી
$!B, \Gamma \fCenter \Delta', !B \lif !C, !C, !A$ની
!!a{proof}~$\pi_1'$ મેળવીએ
(ડાબે~$!B$ અને જમણે~$!C$ ઉમેરીએ).
એ જ રીતે~$\pi_3$ પર
\olref[seq][adm]{prop:weak-G3c-adm} લગાવી
$!A, !B, \Gamma \fCenter \Delta', !B \lif !C, !C$ની
!!a{proof}~$\pi_3'$ મેળવીએ
(જમણે~$!B \lif !C$ ઉમેરીને).
આગમન-પૂર્વધારણા નીચેની સાબિતી પર લાગુ પડે છે:
\[
\AxiomC{}
\RightLabel{$\pi_1'$}
\Deduce$!B, \Gamma \fCenter \Delta', !B \lif !C, !C, !A$
\AxiomC{}
\RightLabel{$\pi_3'$}
\Deduce$!A, !B, \Gamma \fCenter \Delta', !B \lif !C, !C$
\RightLabel{\CutCS}
\BinaryInf$!B, \Gamma \fCenter \Delta', !B \lif !C, !C$
\DisplayProof
\]
કટ !!{formula}~$!A$ છે. આ !!{proof}નો
કટ-ક્રમ~$\pi$ જેટલો જ છે,
પણ કટની !!{height} ઓછી છે.
આથી \Log{G3c}માં \CutCS{} વિના
$!B, \Gamma \fCenter \Delta', !B \lif !C, !C$ની
!!a{proof}~$\pi_4$ મળે છે.
તેના પર~$\RightR{\lif}$ લગાવતાં
$\Gamma \fCenter \Delta', !B \lif !C, !B \lif !C$ની
!!a{proof} મળે છે:
\[
\AxiomC{}
\RightLabel{$\pi_4$}
\Deduce$!B, \Gamma \fCenter \Delta', !B \lif !C, !C$
\RightLabel{$\RightR{\lif}$}
\UnaryInf$\Gamma \fCenter \Delta', !B \lif !C, !B \lif !C$
\DisplayProof
\]
$\Gamma \Sequent \Delta$,
એટલે~$\Gamma \fCenter \Delta', !B \lif !C$ની
!!a{proof}~$\pi'$ મેળવવા
\olref[seq][inv]{prop:G3c-cont-adm}
(સંકોચનની સ્વીકાર્યતા) લગાવો.

હવે માનો કે~$\pi$નો અંત આ રીતે થાય છે:
\[
\AxiomC{}
\RightLabel{$\pi_1$}
\Deduce$\lexists[x][!B(x)], \Gamma' \fCenter \Delta, !A$
\AxiomC{}
\RightLabel{$\pi_3$}
\Deduce$!A, !B(c), \Gamma' \fCenter \Delta$
\RightLabel{$\LeftR{\lexists}$}
\UnaryInf$!A, \lexists[x][!B(x)], \Gamma' \fCenter \Delta$
\RightSubproofLabel{$\pi_2$}
\RightLabel{\CutCS}
\BinaryInf$\Gamma' \fCenter \Delta$
\DisplayProof
\]
ફરી આગમન-પૂર્વધારણા આના પર લગાવીએ:
\[
\AxiomC{}
\RightLabel{$\pi_1[!B(c) \Sequent {}]$}
\Deduce$\lexists[x][!B(x)], !B(c), \Gamma' \fCenter \Delta, !A$
\AxiomC{}
\LeftLabel{$\pi_3[\lexists[x][!B(x)] \Sequent {}]$}
\Deduce$!A, \lexists[x][!B(x)], !B(c), \Gamma' \fCenter \Delta$
\RightLabel{\CutCS}
\BinaryInf$\lexists[x][!B(x)], !B(c), \Gamma' \fCenter \Delta$
\RightLabel{\LeftR{\lexists}}
\UnaryInf$\lexists[x][!B(x)], \lexists[x][!B(x)], \Gamma' \fCenter \Delta$
\DisplayProof
\]
ધ્યાન આપો કે સ્વચલની શરત પૂરી થાય છે:
~$\pi_2$માં \LeftR{\lexists}ની સ્વચલ-શરત
ખાતરી આપે છે કે~$c$ એ
$\lexists[x][!B(x)]$, $\Gamma'$
કે~$\Delta$માં આવતો નથી.
અંતે \olref[seq][inv]{prop:G3c-cont-adm} લગાવો.


\paragraph{કિસ્સો D:}
$\pi_1$ અને $\pi_2$માંથી ઓછામાં ઓછી એક
એવા બે-આધારવિધાનવાળા નિયમ~$R$થી પૂરી થાય છે
જેમાં~$!A$ મુખ્ય નથી.
છ કિસ્સા છે. અહીં~$\pi_2$નું છેલ્લું
અનુમાન~$\RightR{\land}$ હોય
અને તેમાં~$!A$ મુખ્ય ન હોય તે કિસ્સો
જોઈએ; બીજા સમાન છે.

માનો કે~$\pi$નો અંત આ રીતે થાય છે:
\[
\AxiomC{}
\RightLabel{$\pi_1$}
\Deduce$\Gamma \fCenter \Delta', !B \land !C, !A$
\AxiomC{}
\RightLabel{$\pi_3$}
\Deduce$!A, \Gamma \fCenter \Delta', !B$
\AxiomC{}
\RightLabel{$\pi_4$}
\Deduce$!A, \Gamma \fCenter \Delta', !C$
\RightLabel{$\RightR{\land}$}
\BinaryInf$!A, \Gamma \fCenter \Delta', !B \land !C$
\RightSubproofLabel{$\pi_2$}
\RightLabel{\CutCS}
\BinaryInf$\Gamma \fCenter \Delta$
\DisplayProof
\]
આગમન-પૂર્વધારણા નીચેની !!{proof}s
પર લાગુ પડે છે:
\[
\AxiomC{}
\RightLabel{$\pi_1'$}
\Deduce$\Gamma \fCenter \Delta', !B \land !C, !B, !A$
\AxiomC{}
\RightLabel{$\pi_3'$}
\Deduce$!A, \Gamma \fCenter \Delta', !B \land !C, !B$
\RightLabel{\CutCS}
\BinaryInf$\Gamma \fCenter \Delta', !B \land !C, !B$
\DisplayProof
\]
અને આના પર પણ:
\[
\AxiomC{}
\RightLabel{$\pi_1''$}
\Deduce$\Gamma \fCenter \Delta', !B \land !C, !C, !A$
\AxiomC{}
\RightLabel{$\pi_4'$}
\Deduce$!A, \Gamma \fCenter \Delta', !B \land !C, !C$
\RightLabel{\CutCS}
\BinaryInf$\Gamma \fCenter \Delta', !B \land !C, !C$
\DisplayProof
\]
\sourcecorrection{OLMD-202}{મૂળ બીજા કટ-વૃક્ષમાં $\Delta$ પછી ફરી $!B\land !C$ ઉમેરે છે, જોકે $\Delta=\Delta',!B\land !C$; અહીં $\Delta'$ રાખવાથી બંને આધારવિધાન અને નિષ્કર્ષ મેળ ખાય છે.}
અહીં \Log{G3c}ની !!{proof}s
$\pi_1'$ અને $\pi_1''$ એ~$\pi_1$માંથી
!!{height} જાળવતા શિથિલીકરણથી
(\olref[seq][adm]{prop:weak-G3c-adm}) મળે છે
(એક વાર જમણે~$!B$ અને બીજી વાર~$!C$
ઉમેરીને). !!{proof}s
$\pi_3'$ અને~$\pi_4'$ પણ અનુક્રમે
$\pi_3$ અને~$\pi_4$માંથી
!!{height} જાળવતા શિથિલીકરણથી મળે છે
(જમણે~$!B \land !C$ ઉમેરીને).
આમ~$\pi_3$ અને~$\pi_4$ પર પણ
શિથિલીકરણ લાગુ પડે છે.

આગમન-પૂર્વધારણાથી \Log{G3c}માં
$\Gamma \fCenter \Delta', !B \land !C, !B$
અને $\Gamma \fCenter \Delta', !B \land !C, !C$ની
!!{proof}s~$\pi_5$ અને~$\pi_6$ મળે છે.
નિયમ~$\RightR{\land}$ વડે તેમને
$\Gamma \Sequent \Delta', !B \land !C, !B \land !C$ની
!!a{proof}માં જોડીએ:
\[
\AxiomC{}
\RightLabel{$\pi_5$}
\Deduce$\Gamma \fCenter \Delta', !B \land !C, !B$
\AxiomC{}
\RightLabel{$\pi_6$}
\Deduce$\Gamma \fCenter \Delta', !B \land !C, !C$
\RightLabel{\RightR{\land}}
\BinaryInf$\Gamma \fCenter \Delta', !B \land !C, !B \land !C$
\DisplayProof
\]
\sourcecorrection{OLMD-203}{મૂળ સંયુક્ત વૃક્ષના અંતિમ સિક્વન્ટમાં $\Delta$ પછી $!B\land !C$ની બે નકલો વધુ છે; $\Delta'$ સાથે બે નકલો મળે છે અને સંકોચનથી ઇચ્છિત $\Delta$ મળે છે.}
$\Gamma \Sequent \Delta$,
એટલે~$\Gamma \fCenter \Delta', !B \land !C$ની
!!a{proof} માટે
\olref[seq][inv]{prop:G3c-cont-adm}
(સંકોચનની સ્વીકાર્યતા) લગાવો.

\paragraph{કિસ્સો E: કટ ઘટાડવા.}
છેલ્લો કિસ્સો ત્યારે છે જ્યારે
$\pi_1$ અને $\pi_2$ બંને એવા
અનુમાનોથી પૂરી થાય છે જેમાં~$!A$
મુખ્ય છે. ~ $!A$ના દરેક શક્ય
!!{main operator} માટે એક,
એમ છ કિસ્સા છે.

દરેક કિસ્સામાં કટ !!{formula}~$!A$
ધરાવતા \CutCS{}ને~$!A$ના
ઉપ-!!{formula}s પરના એક અથવા વધુ કટમાં
ફેરવીએ છીએ. આ એ અનુમાનોના સક્રિય
!!{formula}s છે જેમાં~$!A$ મુખ્ય
!!{formula} છે. તેથી આ કિસ્સાઓને
કટ અનુમાનનું ``ઘટાડવું'' અથવા
``રૂપાંતર'' પણ કહેવાય છે.

\begin{enumerate}
\item જો $!A \ident \lnot !B$ હોય,
તો~$\pi$નો અંત આ રીતે થાય છે:
\[
\AxiomC{}
\RightLabel{$\pi_1'$}
\Deduce$!B, \Gamma \fCenter \Delta$
\RightLabel{\RightR{\lnot}}
\UnaryInf$\Gamma \fCenter \Delta, \lnot !B$
\LeftSubproofLabel{$\pi_1$}
\AxiomC{}
\RightLabel{$\pi_2'$}
\Deduce$\Gamma \fCenter \Delta, !B$
\RightLabel{\LeftR{\lnot}}
\UnaryInf$\lnot !B, \Gamma \fCenter \Delta$
\RightSubproofLabel{$\pi_2$}
\RightLabel{\CutCS}
\BinaryInf$\Gamma \fCenter \Delta$
\DisplayProof
\]
આગમન-પૂર્વધારણા વડે આમાંથી
\CutCS{} દૂર કરી શકાય છે:
\[
\AxiomC{}
\RightLabel{$\pi_2'$}
\Deduce$\Gamma \fCenter \Delta, !B$
\AxiomC{}
\RightLabel{$\pi_1'$}
\Deduce$!B, \Gamma \fCenter \Delta$
\RightLabel{\CutCS}
\BinaryInf$\Gamma \fCenter \Delta$
\DisplayProof
\]

\item
જો $!A \ident (!B \lor !C)$ હોય,
તો~$\pi$નો અંત આ રીતે થાય છે:
\[
\AxiomC{}
\RightLabel{$\pi_1'$}
\Deduce$\Gamma \fCenter \Delta, !B, !C$
\RightLabel{\RightR{\lor}}
\UnaryInf$\Gamma \fCenter \Delta, !B \lor !C$
\LeftSubproofLabel{$\pi_1$}
\AxiomC{}
\RightLabel{$\pi_2'$}
\Deduce$!B, \Gamma \fCenter \Delta$
\AxiomC{}
\RightLabel{$\pi_2''$}
\Deduce$!C, \Gamma \fCenter \Delta$
\RightLabel{\LeftR{\lor}}
\BinaryInf$!B \lor !C, \Gamma \fCenter \Delta$
\RightSubproofLabel{$\pi_2$}
\RightLabel{\CutCS}
\BinaryInf$\Gamma \fCenter \Delta$
\DisplayProof
\]
\CutCS{}થી પૂરી થતી આ !!{proof} લો:
\[
\AxiomC{}
\RightLabel{$\pi_1'$}
\Deduce$\Gamma \fCenter \Delta, !B, !C$
\AxiomC{}
\RightLabel{$\pi_2'''$}
\Deduce$!C, \Gamma \fCenter \Delta, !B$
\RightLabel{\CutCS}
\BinaryInf$\Gamma \fCenter \Delta, !B$
\DisplayProof
\]
અહીં~$\pi_2'''$ એ~$\pi_2''$માંથી
!!{height} જાળવતા શિથિલીકરણથી મળે છે.
તેનો કટ-ક્રમ $< \cutrank{\pi}$ છે,
કારણ કે $\depth{!C} < \depth{!B \lor !C}$.
\sourcecorrection{OLMD-204}{મૂળમાં અપરિભાષિત $\cutr{\pi}$ છે; આ વિભાગની વ્યાખ્યા પ્રમાણે યોગ્ય સંજ્ઞા $\cutrank{\pi}$ છે.}
તેથી આગમન-પૂર્વધારણાથી
\Log{G3c}માં
$\Gamma \fCenter \Delta, !B$ની
!!a{proof}~$\pi_1''$ મળે છે.
હવે \CutCS{}થી પૂરી થતી આ !!{proof} લો:
\[
\AxiomC{}
\RightLabel{$\pi_1''$}
\Deduce$\Gamma \fCenter \Delta, !B$
\AxiomC{}
\RightLabel{$\pi_2'$}
\Deduce$!B, \Gamma \fCenter \Delta$
\RightLabel{\CutCS}
\BinaryInf$\Gamma \fCenter \Delta$
\DisplayProof
\]
તેનો કટ-ક્રમ
$= \depth{!B} < \depth{!B \lor !C}$ છે.
આથી આગમન-પૂર્વધારણા ફરી લાગુ પડે છે
અને $\Gamma \fCenter \Delta$ની
!!a{proof}~$\pi'$ મળે છે.

\item $!A \ident (!B \land !C)$. કસરત.

\item જો $!A \ident (!B \lif !C)$ હોય,
તો~$\pi$નો અંત આ રીતે થાય છે:
\[
\AxiomC{}
\RightLabel{$\pi_1'$}
\Deduce$!B, \Gamma \fCenter \Delta, !C$
\RightLabel{\RightR{\lif}}
\UnaryInf$\Gamma \fCenter \Delta, !B \lif !C$
\LeftSubproofLabel{$\pi_1$}
\AxiomC{}
\RightLabel{$\pi_2'$}
\Deduce$\Gamma \fCenter \Delta, !B$
\AxiomC{}
\RightLabel{$\pi_2''$}
\Deduce$!C, \Gamma \fCenter \Delta$
\RightLabel{\LeftR{\lif}}
\BinaryInf$!B \lif !C, \Gamma \fCenter \Delta$
\RightSubproofLabel{$\pi_2$}
\RightLabel{\CutCS}
\BinaryInf$\Gamma \fCenter \Delta$
\DisplayProof
\]
\CutCS{}થી પૂરી થતી આ !!{proof} લો:
\[
\AxiomC{}
\RightLabel{$\pi_1'$}
\Deduce$!B, \Gamma \fCenter \Delta, !C$
\AxiomC{}
\RightLabel{$\pi_2''$}
\Deduce$!C, \Gamma \fCenter \Delta$
\RightLabel{\CutCS}
\BinaryInf$!B, \Gamma \fCenter \Delta$
\DisplayProof
\]
તેનો કટ-ક્રમ~$\pi$ કરતાં ઓછો છે.
આગમન-પૂર્વધારણા
$!B, \Gamma \fCenter \Delta$ની
!!a{proof}~$\pi_1''$ આપે છે.
હવે \CutCS{}થી પૂરી થતી આ !!{proof} લો:
\[
\AxiomC{}
\RightLabel{$\pi_2'$}
\Deduce$\Gamma \fCenter \Delta, !B$
\AxiomC{}
\RightLabel{$\pi_1''$}
\Deduce$!B, \Gamma \fCenter \Delta$
\RightLabel{\CutCS}
\BinaryInf$\Gamma \fCenter \Delta$
\DisplayProof
\]
તેનો કટ-ક્રમ પણ~$\pi$ના કટ-ક્રમ કરતાં
ઓછો છે. તેથી આગમન-પૂર્વધારણા
$\Gamma \fCenter \Delta$ની
\Log{G3c}-!!{proof}~$\pi'$ આપે છે.
\item જો $!A \ident \lforall[x][!B(x)]$ હોય,
તો~$\pi$નો અંત આ રીતે થાય છે:
\[
\AxiomC{}
\RightLabel{$\pi_1'$}
\Deduce$\Gamma \fCenter \Delta, !B(c)$
\RightLabel{\RightR{\lforall}}
\UnaryInf$\Gamma \fCenter \Delta, \lforall[x][!B(x)]$
\LeftSubproofLabel{$\pi_1$}
\AxiomC{}
\RightLabel{$\pi_2'$}
\Deduce$!B(t), \lforall[x][!B(x)], \Gamma \fCenter \Delta$
\RightLabel{\LeftR{\lforall}}
\UnaryInf$\lforall[x][!B(x)], \Gamma \fCenter \Delta$
\RightSubproofLabel{$\pi_2$}
\RightLabel{\CutCS}
\BinaryInf$\Gamma \fCenter \Delta$
\DisplayProof
\]
આગમન-પૂર્વધારણા વડે આમાંથી
\CutCS{} દૂર કરી શકાય છે:
\[
\AxiomC{}
\RightLabel{$\pi_1''$}
\Deduce$!B(t), \Gamma \fCenter \Delta, \lforall[x][!B(x)]$
\AxiomC{}
\RightLabel{$\pi_2'$}
\Deduce$!B(t), \lforall[x][!B(x)], \Gamma \fCenter \Delta$
\RightLabel{\CutCS}
\BinaryInf$!B(t), \Gamma \fCenter \Delta$
\DisplayProof
\]
અહીં~$\pi_1''$ એ~$\pi_1$માંથી
(~$\pi_1'$માંથી નહીં!)
!!{height} જાળવતા શિથિલીકરણથી મળે છે.
(આ !!{proof}નો કટ-ક્રમ સમાન છે,
પણ કટની !!{height} ઓછી છે.)
તેથી $!B(t), \Gamma \fCenter \Delta$ની
!!a{proof}~$\pi_2''$ મળે છે.

!!{proof}~$\pi_1'$નો અંતિમ સિક્વન્ટ
$\Gamma \Sequent \Delta, !B(c)$ છે,
જ્યાં~$c$ એ~$\RightR{\lforall}$
અનુમાનનો સ્વચલ છે.
તેથી~$c$ એ~$!B(x)$, $\Gamma$
કે~$\Delta$માં આવતો નથી.
!!{proof}~$\pi$ નિયમિત છે એમ ધારીએ,
એટલે~$c$ માત્ર ત્યાર પછી આવતા
$\RightR{\lforall}$ અનુમાનનો જ સ્વચલ છે
અને તેની ઉપર જ આવે છે;
અર્થાત્, તે માત્ર~$\pi_1'$માં આવે છે
અને~$\pi_1'$ની અંદરના કોઈ અનુમાનનો
સ્વચલ નથી. કારણ કે~$c$ એ~$\pi_2$માં
આવતો નથી, તે~$t$માં પણ આવતો નથી.
\olref[seq][qua]{lem:subst-regular} મુજબ
!!{proof}~$\Subst{\pi_1'}{t}{c}$ એ
$\Gamma \Sequent \Delta, !B(t)$ની
!!a{proof} છે. હવે નીચેની
!!{proof} પર આગમન-પૂર્વધારણા લગાવીએ:
\[
\AxiomC{}
\RightLabel{$\Subst{\pi_1'}{t}{c}$}
\Deduce$\Gamma \fCenter \Delta, !B(t)$
\AxiomC{}
\RightLabel{$\pi_2''$}
\Deduce$!B(t), \Gamma \fCenter \Delta$
\RightLabel{\CutCS}
\BinaryInf$\Gamma \fCenter \Delta$
\DisplayProof
\]
(આગમન-પૂર્વધારણા લાગુ પડે છે,
કારણ કે કટ !!{formula}~$!B(t)$ છે;
એટલે આ !!{proof}નો કટ-ક્રમ~$\pi$
કરતાં ઓછો છે.)
\item $!A \ident \lexists[x][!B(x)]$નો
કિસ્સો કસરત તરીકે છોડીએ છીએ.
\end{enumerate}
\end{proof}

\begin{prob}
\olref{lem:cut-adm-G3c}ની સાબિતીના કિસ્સા~Dનો
એ ઉપકિસ્સો ચકાસો જેમાં~$!A$ એ~$\pi_1$ના
છેલ્લા અનુમાનમાં મુખ્ય નથી અને તે અનુમાનનો
અંત~$\LeftR{\lif}$થી થાય છે.
(ધ્યાન આપો: આ કસરતનો એક ભાગ એ પણ છે કે
શું સાબિત કરવાનું છે તે સ્પષ્ટ કહો,
એટલે આ કિસ્સામાં~$\pi$નું સ્વરૂપ જણાવો.)
\end{prob}

\begin{prob}
\olref{lem:cut-adm-G3c}ની સાબિતીના કિસ્સા~Dનો
એ ઉપકિસ્સો ચકાસો જેમાં~$!A$ એ~$\pi_1$ના
છેલ્લા અનુમાનમાં મુખ્ય નથી અને તે અનુમાનનો
અંત~$\LeftR{\lforall}$થી થાય છે.
\end{prob}

\begin{prob}
\olref{lem:cut-adm-G3c}ની સાબિતીના કિસ્સા~Eનો
એ ઉપકિસ્સો ચકાસો જેમાં~$!A$ એ
$\pi_1$ અને $\pi_2$ બંનેના
છેલ્લા અનુમાનમાં મુખ્ય છે અને
$!A \ident (!B \land !C)$ છે.
\end{prob}

\begin{prob}
\olref{lem:cut-adm-G3c}ની સાબિતીના કિસ્સા~Eનો
એ ઉપકિસ્સો ચકાસો જેમાં~$!A$ એ
$\pi_1$ અને $\pi_2$ બંનેના
છેલ્લા અનુમાનમાં મુખ્ય છે અને
$!A \ident \lexists[x][!B(x)]$ છે.
\end{prob}

ધ્યાન આપો કે કિસ્સા~Eમાં આગમન-પૂર્વધારણા
જે~\CutCS{}થી પૂરી થતી !!{proof} પર
લગાડીએ તેની કટની !!{height} મૂળ
સાબિતીની કટની !!{height} કરતાં વધારે
હોઈ શકે. દાખલા તરીકે,
$!A \ident (!B \lif !C)$ હોય ત્યારે
~$\pi$ને બે \CutCS{} અનુમાનો ધરાવતી
!!a{proof}માં ફેરવ્યું:
\[
\AxiomC{}
\RightLabel{$\pi_2'$}
\Deduce$\Gamma \fCenter \Delta, !B$
\AxiomC{}
\RightLabel{$\pi_1'$}
\Deduce$!B, \Gamma \fCenter \Delta, !C$
\AxiomC{}
\RightLabel{$\pi_2''$}
\Deduce$!C, \Gamma \fCenter \Delta$
\RightLabel{\CutCS}
\BinaryInf$!B, \Gamma \fCenter \Delta$
\RightSubproofLabel{$\pi_1''$}
\RightLabel{\CutCS}
\BinaryInf$\Gamma \fCenter \Delta$
\DisplayProof
\]
\sourcecorrection{OLMD-205}{મૂળ અંતિમ વૃક્ષમાં બીજા કટ પછી પણ પૂર્વાંગમાં $!B$ રહી જાય છે; $!B$ પર કટ કરતાં અંતિમ સિક્વન્ટ $\Gamma\Sequent\Delta$ છે.}
ઉપરનો \CutCS{} એ~$\pi_1'$ અને $\pi_2''$
વચ્ચે છે; તેનો કટ-ક્રમ અને કટની
!!{height} બંને~$\pi$ કરતાં ઓછાં છે.
પણ આગમન-પૂર્વધારણાથી મળેલી
!!{proof}~$\pi_1''$ની કટ-!!{height}
હવે~$\pi$ કરતાં ઓછી હોવી જરૂરી નથી.
તેમ છતાં, નીચેના \CutCS{}નું
કટ !!{formula}~$!B$ હોવાથી તેનો
કટ \emph{ક્રમ}~$\pi$ના કટ-ક્રમ કરતાં
ઓછો છે.

\olref{lem:cut-adm-G3c} બતાવે છે કે
!!a{proof}માંથી સૌથી ઉપરનું એક
\CutCS{} અનુમાન દૂર કરી શકાય.
સૌથી ઉપરના~\CutCS{}માં પૂરી થતી
ઉપ-!!{proof}s પર તેને વારંવાર
લગાવી કોઈ પણ !!a{proof}માંથી
ગમે તેટલાં \CutCS{} અનુમાનો
દૂર કરી શકાય તે બતાવીશું.

\begin{thm}\ollabel{thm:cut-elim-G3c-topmost}
$\Log{G3c} + \CutCS$ની કોઈ પણ
!!{proof}~$\pi$ને~\Log{G3c}ની
!!a{proof}માં ફેરવી શકાય છે.
\end{thm}

\begin{proof}
  ~$\pi$માંનાં \CutCS{} અનુમાનોની
  સંખ્યા~$n$ પર આગમન કરીએ.
  જો $n=0$, તો કશું કરવાનું નથી:
  $\pi$ પહેલેથી~\Log{G3c}ની
  !!a{proof} છે. અન્યથા, સૌથી ઉપરના
  \CutCS{} અનુમાનમાં પૂરી થતી
  ઉપ-!!{proof}~$\pi'$ લો.
  તેમાં એક જ \CutCS{} નિયમ છે,
  અને તે છેલ્લું અનુમાન છે.
  \olref{lem:cut-adm-G3c} લગાવી
  તેને \CutCS{} વિનાની
  સાબિતી~$\pi''$માં ફેરવો,
  અને~$\pi'$ ઉપ-!!{proof}ને
  $\pi''$થી બદલો.
  મળેલી !!a{proof}માં
  $n-1$ \CutCS{} અનુમાનો રહે છે,
  તેથી આગમન-પૂર્વધારણા લાગુ પડે છે.
\end{proof}

\end{document}
